% ลำดับขั้นย่อยภายในชั้นจัดอันดับ — แทรกก่อนหัวข้อการรวมสัญญาณในบทที่ 4

\subsection{ลำดับการทำงานภายในชั้นจัดอันดับ}

ชั้นจัดอันดับรับผู้เข้าชิง 100 รายการจากขั้นค้นคืน แล้วทำงาน 4 ขั้นย่อยตามลำดับ
ดังตารางที่ \ref{tab:rankorder}

\begin{table}[H]
\centering
\caption{ลำดับขั้นย่อยภายในชั้นจัดอันดับ พร้อมจำนวนที่เหลือจากตัวอย่างการรันจริง}
\label{tab:rankorder}
\small
\begin{tabular}{|c|p{7.2cm}|c|}
\hline
\textbf{ขั้นย่อย} & \textbf{การทำงาน} & \textbf{จำนวนที่เหลือ} \\
\hline
1 & ตัดทิ้งตามเงื่อนไขที่ผู้ใช้ระบุเอง เช่น ราคาเกินงบ หรือแบรนด์ที่ผู้ใช้ปฏิเสธ & 100 \\ \hline
2 & คัดด้วยเกณฑ์สัมบูรณ์ 0.55 โดยใช้คะแนนความใกล้เชิงความหมายเพียงอย่างเดียว & 73 \\ \hline
3 & \textbf{จัดอันดับด้วย 5 สัญญาณ รวมถึงการนับคำซ้อน} & 73 \\ \hline
4 & หยิบ 3 รายการแรก โดยเลี่ยงการหยิบจากร้านค้าเดียวกันซ้ำ & 3 \\
\hline
\end{tabular}
\end{table}

\textbf{ประเด็นที่ต้องเข้าใจให้ตรงคือการนับคำซ้อนไม่ได้ทำหน้าที่คัดสินค้าทิ้ง}
สัญญาณนี้ทำงานที่ขั้นย่อยที่ 3 ซึ่งอยู่หลังการคัดด้วยเกณฑ์สัมบูรณ์แล้ว
หน้าที่ของมันคือเรียงลำดับสินค้าที่ผ่านเกณฑ์มาแล้วเท่านั้น
สินค้าที่ได้คะแนนการนับคำซ้อนเป็นศูนย์จึงยังอยู่ในกอง เพียงแต่ถูกจัดไปอยู่อันดับท้าย

การแยกหน้าที่เช่นนี้จำเป็น เพราะการคัดว่าดีพอหรือยังต้องใช้เกณฑ์สัมบูรณ์
ส่วนการจัดอันดับเป็นการเปรียบเทียบกันเองภายในกอง ถ้านำการนับคำซ้อนไปใช้เป็นเกณฑ์คัด
สินค้าที่คลังเรียกด้วยคำอื่นจะถูกตัดทิ้งทั้งที่ตรงกับความต้องการ
ซึ่งเป็นข้อผิดพลาดที่พบจริงและรายงานไว้ในหัวข้อผลเชิงลบ

ข้อความที่นำมานับคำซ้อนคือคำค้นหลังผ่านขั้นแปลภาษาแล้ว ไม่ใช่ข้อความที่ผู้ใช้พิมพ์โดยตรง
ฝั่งภาษาไทยจึงได้รับผลกระทบจากความคลาดเคลื่อนของขั้นแปลอย่างหลีกเลี่ยงไม่ได้

นอกจากนี้ ชั้นจัดอันดับถูกเรียกใช้ทุกครั้งที่มีการค้น ไม่ใช่เพียงครั้งเดียวต่อหนึ่งคำถาม
เมื่อการค้นซ้ำทำงาน ผู้เข้าชิงชุดใหม่จะถูกจัดอันดับใหม่ทั้งหมด
และชั้นนี้เปิดช่องให้เปลี่ยนน้ำหนักของการนับคำซ้อนขณะทำงานได้
ซึ่งเป็นช่องทางที่ใช้ทดลองให้เอเจนต์ปรับน้ำหนักเอง และพบว่าไม่ได้ผล
